\documentclass[12pt]{article}
\usepackage[utf8]{inputenc}
\usepackage[T1]{fontenc}
\usepackage[margin=1in]{geometry}
\usepackage{setspace}
\usepackage{xcolor}
\usepackage[colorlinks=true, urlcolor=blue!60!black, linkcolor=blue!60!black]{hyperref}
\setlength{\parindent}{2em}
\setlength{\parskip}{1.25em}
\renewcommand{\baselinestretch}{1.0}
\begin{document}
\begin{center}
{\LARGE \textit{TransUnion LLC v. Ramirez}, 594 U.S. 413 (2021)\par}
\vspace{0.8em}
{\large Case Brief\par}
\end{center}
\vspace{1em}

\section{Memory Jogger}

\textbf{``No concrete harm, no standing''---the unsent defamatory letter in the desk drawer.} When credit reporting agency TransUnion falsely flagged 8,185 consumers as potential terrorists on a federal watchlist in violation of the Fair Credit Reporting Act, the Supreme Court held 5--4 that only the 1,853 class members whose misleading credit reports were actually disseminated to third-party businesses suffered a ``concrete'' injury in fact under Article III. The other 6,332 class members---whose inaccurate alerts sat unshared in internal databases---lacked standing to sue for damages in federal court because ``an injury in law is not an injury in fact,'' and an unmaterialized risk of future harm cannot by itself confer standing in a suit for retrospective damages (\href{https://www.supremecourt.gov/opinions/20pdf/20-297_4g25.pdf}{Slip Opinion}; \href{https://supreme.justia.com/cases/federal/us/594/20-297/}{Justia}).

\section{Detailed Case Facts}

TransUnion LLC is one of the ``Big Three'' consumer credit reporting agencies in the United States, alongside Equifax and Experian. Under the Fair Credit Reporting Act (FCRA), 15 U.S.C.~\S~1681 \textit{et seq.}, consumer reporting agencies must: (1) ``follow reasonable procedures to assure maximum possible accuracy'' in consumer reports (\S~1681e(b)); (2) disclose to a consumer, upon request, ``[a]ll information in the consumer's file at the time of the request'' (\S~1681g(a)(1)); and (3) provide a Consumer Financial Protection Bureau (CFPB) ``summary of rights'' with each written disclosure (\S~1681g(c)(2)). Section 1681n(a) creates a private cause of action for willful violations, authorizing actual damages or statutory damages between \$100 and \$1,000, plus punitive damages and attorney's fees (\href{https://www.law.cornell.edu/supremecourt/text/20-297}{Cornell LII}).

Beginning in 2002, TransUnion marketed an add-on product called ``OFAC Name Screen Alert.'' The U.S.~Treasury Department's Office of Foreign Assets Control (OFAC) maintains a list of ``Specially Designated Nationals''---terrorists, drug traffickers, and other serious criminals with whom Americans are generally barred from transacting business. When a business subscribed to Name Screen, TransUnion used third-party software to compare only the consumer's first and last name against the OFAC list, without checking birth dates, middle initials, Social Security numbers, or other identifiers. Predictably, this rudimentary matching generated thousands of false positives for law-abiding Americans (\href{https://supreme.justia.com/cases/federal/us/594/20-297/}{Justia}). Despite a prior 2005 lawsuit and jury verdict against TransUnion over the same product (\textit{Cortez v. Trans Union, LLC}, 617 F.3d 688 (3d Cir. 2010)), TransUnion made minimal changes, shifting only from labeling consumers a ``match'' to a ``potential match'' and requiring exact first-and-last-name matches (\href{https://supreme.justia.com/cases/federal/us/594/20-297/}{Justia}).

On February 27, 2011, Sergio Ramirez visited a Nissan dealership in Dublin, California, with his wife and father-in-law to purchase a Nissan Maxima. After they negotiated a price, the dealership ran a credit check. Ramirez's TransUnion credit report displayed the alert: ``***OFAC ADVISOR ALERT - INPUT NAME MATCHES NAME ON THE OFAC DATABASE.'' A Nissan salesperson refused to sell Ramirez the car because his name appeared on a ``terrorist list,'' forcing his wife to purchase the vehicle in her own name (\href{https://www.law.cornell.edu/supremecourt/text/20-297}{Cornell LII}).

The next day, Ramirez contacted TransUnion and requested a copy of his credit file. TransUnion sent him two separate mailings: first, a credit file and CFPB summary of rights that omitted any mention of the OFAC alert; and second, a ``courtesy'' letter the following day informing him that his name was a potential match on the OFAC list, without disclosing that the alert appeared on his credit report and without enclosing another CFPB summary of rights. Alarmed, Ramirez consulted an attorney and canceled a planned trip to Mexico. TransUnion later removed the OFAC alert from his file (\href{https://www.supremecourt.gov/opinions/20pdf/20-297_4g25.pdf}{Slip Opinion}; \href{https://supreme.justia.com/cases/federal/us/594/20-297/}{Justia}).

\section{Procedural History}

In February 2012, Ramirez sued TransUnion in the U.S.~District Court for the Northern District of California, alleging three willful FCRA violations: (1) failure to follow reasonable procedures to assure maximum possible accuracy (\S~1681e(b)); (2) failure to disclose all information in the consumer's file upon request (\S~1681g(a)(1)); and (3) failure to include a summary of rights with each written disclosure (\S~1681g(c)(2)). Ramirez sought statutory and punitive damages on behalf of a putative nationwide class of consumers who received a similar OFAC mailing from TransUnion between January 1, 2011, and July 26, 2011 (\href{https://supreme.justia.com/cases/federal/us/594/20-297/}{Justia}).

The district court certified the class over TransUnion's objection (301 F.R.D. 408 (N.D. Cal. 2014)). Before trial, the parties stipulated that the class comprised 8,185 individuals (including Ramirez), and that TransUnion had disseminated credit reports containing OFAC alerts to third-party businesses for only 1,853 class members during the seven-month class period; the internal credit files of the remaining 6,332 class members were not provided to third parties during that period. Following the Supreme Court's decision in \textit{Spokeo, Inc. v. Robins}, 578 U.S. 330 (2016), the district court held that all 8,185 class members possessed Article III standing on all three claims (2016 WL 6070490, at *5 (N.D. Cal. Oct. 17, 2016)). After a six-day trial at which Ramirez testified about his own experience but presented no evidence of other class members' experiences, the jury returned a verdict for the class on all three claims, awarding each class member \$984.22 in statutory damages (roughly \$8 million total) and \$6,353.08 in punitive damages (roughly \$52 million total), for a combined verdict exceeding \$60 million (\href{https://www.law.cornell.edu/supremecourt/text/20-297}{Cornell LII}; \href{https://harvardlawreview.org/print/vol-135/transunion-v-ramirez/}{Harvard Law Review}).

On appeal, a divided panel of the U.S.~Court of Appeals for the Ninth Circuit affirmed in relevant part (951 F.3d 1008 (9th Cir. 2020)). Writing for the panel (joined by Judge Fletcher), Judge Murguia held that every class member must establish Article III standing at the final judgment stage, but concluded that all 8,185 class members had standing on all three claims because TransUnion's reckless practices exposed every member to a real risk of harm to concrete privacy, reputational, and informational interests. The panel also upheld Rule 23 typicality, but reduced the punitive damages award to \$3,936.88 per class member, lowering the total judgment to approximately \$40 million. Judge McKeown concurred in part and dissented in part, concluding that only the 1,853 class members whose reports were actually disseminated had standing on the reasonable-procedures claim and that only Ramirez had standing on the two mailing claims (\href{https://supreme.justia.com/cases/federal/us/594/20-297/}{Justia}; \href{https://harvardlawreview.org/print/vol-135/transunion-v-ramirez/}{Harvard Law Review}).

The Supreme Court granted certiorari on December 16, 2020, limited to the first question presented concerning Article III standing and Rule 23 typicality (declining review of the punitive damages question), heard oral argument on March 30, 2021, and decided the case on June 25, 2021 (\href{https://www.scotusblog.com/cases/transunion-llc-v-ramirez/}{SCOTUSblog}).

\section{Judicial Votes}

\textbf{Vote:} 5--4, reversing and remanding the judgment of the Ninth Circuit (\href{https://www.scotusblog.com/cases/transunion-llc-v-ramirez/}{SCOTUSblog}; \href{https://supreme.justia.com/cases/federal/us/594/20-297/}{Justia}).

\begin{itemize}
\item \textbf{Majority Opinion (5):} Justice Kavanaugh delivered the opinion of the Court, joined in full by Chief Justice Roberts and Justices Alito, Gorsuch, and Barrett.
\item \textbf{Principal Dissent (4):} Justice Thomas filed a dissenting opinion, joined in full by Justices Breyer, Sotomayor, and Kagan. (This produced an unusual ideological alignment uniting Justice Thomas with the Court's three liberal Justices.)
\item \textbf{Second Dissent (3):} Justice Kagan filed a separate dissenting opinion, joined by Justices Breyer and Sotomayor, agreeing with Justice Thomas's dissent subject to one qualification regarding \textit{Spokeo}.
\end{itemize}

\section{Holding}

\textbf{Judgment:} Reversed and remanded. The Court held that: (1) on the \S~1681e(b) reasonable-procedures claim, only the 1,853 class members whose credit reports containing misleading OFAC alerts were disseminated to third-party businesses had Article III standing, whereas the remaining 6,332 class members lacked standing; and (2) on the \S~1681g(a)(1) disclosure and \S~1681g(c)(2) summary-of-rights claims, none of the 8,185 class members other than named plaintiff Sergio Ramirez had standing. Because its Article III ruling was dispositive, the Court declined to decide whether Ramirez's claims satisfied Rule 23 typicality and left class certification on remand to the Ninth Circuit (\href{https://www.law.cornell.edu/supremecourt/text/20-297}{Cornell LII}).

\textbf{Rule Announced:} Only plaintiffs who suffer a ``concrete harm'' from a defendant's statutory violation have Article III standing to seek damages against a private defendant in federal court. Every class member must demonstrate Article III standing in order to recover individual damages. To be ``concrete,'' an asserted intangible harm must bear a ``close relationship''---though not an ``exact duplicate''---to a harm traditionally recognized as providing a basis for a lawsuit in American courts (such as defamation, disclosure of private information, or intrusion upon seclusion). An ``injury in law is not an injury in fact,'' and in a suit for retrospective damages, the mere risk of future harm, standing alone without materialization or an independent present injury caused by exposure to the risk, does not constitute concrete harm. Finally, an asserted ``informational injury'' without ``downstream consequences'' or ``adverse effects'' fails to satisfy Article III (\href{https://supreme.justia.com/cases/federal/us/594/20-297/}{Justia}).

\textbf{Holding vs.\ Dicta:}
\begin{itemize}
\item \textit{Core Holdings:} (a) Dissemination of a credit report falsely labeling a consumer a ``potential match'' to the OFAC terrorist list bears a close relationship to defamation and confers Article III standing; (b) retention of inaccurate information in an internal database without third-party dissemination does not constitute concrete harm; (c) an unmaterialized risk of future harm cannot by itself confer Article III standing in a suit for damages; (d) formatting defects in required disclosures that cause no downstream adverse effects do not confer standing; and (e) every class member must have Article III standing to recover individual damages at final judgment.
\item \textit{Dicta and Reserved Questions:} The Court expressly reserved (1) whether every class member must demonstrate standing \textit{prior to class certification} (594 U.S. at 431 n.4; slip op. at 15 n.4); (2) whether a plaintiff's knowledge of exposure to a risk of future harm could itself cause an independent emotional or psychological injury sufficient for Article III standing, such as by analogy to intentional infliction of emotional distress (slip op. at 21 n.7); and (3) whether TransUnion violated \S~1681e(b) on the merits with respect to undisseminated files (slip op. at 16 n.5). Justice Kavanaugh's discussion of the Maine-versus-Hawaii environmental hypothetical and his broad separation-of-powers rationale linking private statutory-damages suits to infringement of the Executive Branch's Article II enforcement discretion also sweep well beyond the FCRA claims adjudicated (\href{https://supreme.justia.com/cases/federal/us/594/20-297/}{Justia}; \href{https://harvardlawreview.org/print/vol-135/transunion-v-ramirez/}{Harvard Law Review}).
\end{itemize}

\section{Analysis of Opinions}

\subsection{Majority Opinion (Justice Kavanaugh)}

Justice Kavanaugh anchored the Court's opinion in separation of powers. Quoting \textit{Raines v. Byrd}, 521 U.S. 811, 820 (1997), and Justice Scalia's famous question---``What's it to you?''---the majority emphasized that Article III's ``Cases'' and ``Controversies'' requirement ensures federal courts resolve only real disputes affecting individuals rather than exercising a ``roving commission'' over the political branches or private entities (\href{https://www.law.cornell.edu/supremecourt/text/20-297}{Cornell LII}). Drawing on \textit{Spokeo, Inc. v. Robins}, 578 U.S. 330 (2016), Kavanaugh explained that concreteness is measured by history and tradition: tangible physical and monetary harms readily qualify, as do intangible harms with a ``close relationship'' to harms traditionally actionable in American courts (such as reputational harm, disclosure of private information, intrusion upon seclusion, and constitutional harms). While Congress's judgment is ``instructive'' and Congress may ``elevate'' real-world harms, Congress ``may not simply enact an injury into existence.'' Allowing Congress to authorize uninjured private plaintiffs to recover statutory damages for any legal violation would not only contravene Article III but also usurp the Executive Branch's Article II discretion over how aggressively to enforce federal law (\href{https://supreme.justia.com/cases/federal/us/594/20-297/}{Justia}).

Applying these principles, the majority divided the class on the \S~1681e(b) reasonable-procedures claim. For the 1,853 class members whose reports were sent to third-party businesses, publication of an alert labeling them potential terrorists or drug traffickers bore a sufficiently ``close relationship'' to defamation---even though TransUnion argued ``potential match'' was merely misleading rather than literally false---because \textit{Spokeo} does not demand an ``exact duplicate.'' For the remaining 6,332 class members, however, the absence of third-party dissemination was fatal. Publication is ``essential to liability'' in defamation; inaccurate information sitting in an internal company database is legally equivalent to a defamatory letter stored in a desk drawer. The majority rejected the plaintiffs' forfeited ``internal publication'' theory (sharing files with internal employees or print vendors) for lack of common-law pedigree and lack of evidence that anyone actually read the alerts. It likewise rejected the ``risk of future harm'' theory: while an imminent, substantial risk of future harm can support forward-looking \textit{injunctive} relief under \textit{Clapper v. Amnesty International USA}, 568 U.S. 398 (2013), an unmaterialized risk---like a reckless driver swerving behind a commuter who arrives home safely---cannot support \textit{retrospective damages} unless exposure to the risk itself caused a separate concrete injury. Moreover, the plaintiffs failed to prove a serious likelihood of dissemination or that the 6,332 class members even knew of the alerts (\href{https://www.supremecourt.gov/opinions/20pdf/20-297_4g25.pdf}{Slip Opinion}; \href{https://supreme.justia.com/cases/federal/us/594/20-297/}{Justia}).

Turning to the \S~1681g disclosure and summary-of-rights claims, the majority held that none of the class members apart from Ramirez (whose individual standing TransUnion did not meaningfully contest) proved concrete harm. There was no trial evidence that any other class member even opened the dual mailings, felt confused or distressed, or would have corrected their files had the mailings been formatted differently. Rejecting the United States's amicus theory of ``informational injury'' under \textit{FEC v. Akins}, 524 U.S. 11 (1998), and \textit{Public Citizen v. Department of Justice}, 491 U.S. 440 (1989), the Court distinguished those cases as involving public-disclosure ``sunshine'' laws and emphasized that an asserted informational injury lacking ``downstream consequences'' or ``adverse effects'' does not satisfy Article III (\href{https://www.law.cornell.edu/supremecourt/text/20-297}{Cornell LII}).

\subsection{Principal Dissent (Justice Thomas)}

Justice Thomas---joined by Justices Breyer, Sotomayor, and Kagan---mounted a sweeping originalist and historical critique. Tracing Founding-era law from Blackstone and \textit{Entick v. Carrington} (1765) to Justice Story's circuit decision in \textit{Whittemore v. Cutter} (1813) and the First Congress's 1790 Copyright Act, Thomas distinguished between \textit{private rights} (duties owed individually to a person) and \textit{public rights} (duties owed broadly to the community). Historically, a plaintiff suing to vindicate a public right had to show both legal injury (\textit{injuria}) and actual damage (\textit{damnum}), whereas a plaintiff suing for violation of a private right---whether rooted in common law or newly created by statute---needed only to establish the violation of the right itself (\href{https://supreme.justia.com/cases/federal/us/594/20-297/}{Justia}).

Thomas demonstrated that the modern ``injury in fact'' phrase did not appear until \textit{Association of Data Processing Service Organizations, Inc. v. Camp}, 397 U.S. 150 (1970)---180 years after Article III's ratification---as a statutory gloss on the Administrative Procedure Act meant to \textit{broaden} access to courts. Because the three FCRA provisions created duties owed specifically to individual consumers, the jury's verdict finding willful violations of each class member's private statutory rights amply established an Article III case. Even under the majority's own framework, Thomas argued, the class had standing: TransUnion withheld valuable OFAC credit reports that it charged clients extra to receive; nearly 25\% of the class had false alerts sent to creditors in just seven months, showing a substantial risk of harm under \textit{Spokeo}; TransUnion transmitted the files to third-party print vendors, which sufficed as publication at common law; and common sense confirms that receiving a letter flagging you as a potential terrorist is harmful. Finally, in a prescient footnote (n.9), Thomas warned that the majority's holding could be a ``pyrrhic victory'' for defendants: because state courts are not bound by Article III, stripping federal jurisdiction may leave state courts as the exclusive, non-removable forum for federal statutory class actions (\href{https://supreme.justia.com/cases/federal/us/594/20-297/}{Justia}; \href{https://harvardlawreview.org/print/vol-135/transunion-v-ramirez/}{Harvard Law Review}).

\subsection{Second Dissent (Justice Kagan)}

Justice Kagan---joined by Justices Breyer and Sotomayor---wrote separately to emphasize institutional competence and judicial restraint. She accused the majority of inverting standing doctrine from ``a doctrine of judicial modesty into a tool of judicial aggrandizement,'' overriding Congress's judgment for the first time to bar a specific class of statutorily authorized plaintiffs from federal court. Kagan mocked the majority's factual skepticism, asking why it is ``speculative'' that a company whose business is selling credit reports will sell a credit report, or why consumers who specifically requested their credit files would not open the envelopes or be confused by split mailings (\href{https://www.law.cornell.edu/supremecourt/text/20-297}{Cornell LII}).

Kagan parted ways with Justice Thomas on a single doctrinal point: adhering to \textit{Spokeo}, she maintained that Article III still requires a concrete injury or risk of real harm even for statutory violations. In practice, however, she argued that courts should defer heavily to Congress, overriding a legislative authorization to sue ``when but only when Congress could not reasonably have thought that a suit will contribute to compensating or preventing the harm at issue''---a standard that would yield the same result as Thomas's private-rights rule in all but ``highly unusual cases'' (\href{https://supreme.justia.com/cases/federal/us/594/20-297/}{Justia}; \href{https://www.scotusblog.com/2021/06/court-limits-standing-in-credit-reporting-lawsuit/}{SCOTUSblog}).

\section{Examples: Future Applications}

\begin{enumerate}
\item \textbf{Same-Side (No Standing): Unbreached Retention of Data Under a Privacy Statute.} A federal statute requires cable providers to destroy former subscribers' personal data within one year of cancellation and authorizes \$500 in statutory damages per violation. A former subscriber sues in federal court after discovering the company retained his billing address for two years, without alleging any breach, sale, or third-party disclosure. Under \textit{TransUnion}, the plaintiff lacks Article III standing: mere internal retention of data without disclosure lacks a close common-law analogue, and an unmaterialized risk of future breach cannot support a retrospective damages claim.

\item \textbf{Same-Side (No Standing): Technical Formatting Defect in a Debt-Collection Notice.} A debt collector sends a consumer a dunning letter that contains all required debt-validation disclosures under the Fair Debt Collection Practices Act, but prints the consumer's dispute rights in 8-point font on the reverse side rather than the statutorily mandated 12-point font on the front. The consumer never alleges that she failed to read the notice, was confused, or missed a dispute deadline, but sues for statutory damages. Under \textit{TransUnion}, this is a ``bare procedural violation'' and an informational defect lacking ``downstream consequences'' or ``adverse effects,'' so the plaintiff lacks Article III standing.

\item \textbf{Opposite-Side (Standing Established): Unauthorized Disclosure of Viewing History to Ad Networks.} A news website embeds a third-party tracking pixel that automatically transmits subscribers' personally identifiable video-viewing histories to a social media advertising platform without consent, violating the Video Privacy Protection Act. A subscriber sues for statutory damages. Under \textit{TransUnion}, the plaintiff has Article III standing: unlike an unsent letter in a desk drawer, the website actually disclosed private information to a third party, bearing a close relationship to the traditional privacy tort of public disclosure of private facts or intrusion upon seclusion (see \href{https://media.cadc.uscourts.gov/opinions/docs/2025/08/24-7022-2129694.pdf}{\textit{Pileggi v. Washington Newspaper Publishing Co.}}, D.C. Cir. 2025).

\item \textbf{Opposite-Side (Standing Established): Prospective Injunctive Relief Against Imminent Cyber-Exposure.} Hackers breach a healthcare network's unencrypted server containing patients' Social Security numbers and medical diagnoses and post a sample batch for sale on a dark-web forum while threatening to release the rest within 30 days. Patients whose records are in the stolen database sue under a federal data-security statute seeking an \textit{injunction} compelling the hospital to fund immediate credit-monitoring and security remediation. Under \textit{TransUnion}, because the plaintiffs seek forward-looking injunctive relief rather than retrospective damages and face a sufficiently imminent and substantial risk of harm under \textit{Clapper}, they have Article III standing.

\item \textbf{Fence-Sitter (Close Call): Automated Transmission of Sensitive Debt Data to a Third-Party Print Vendor.} A debt collector transmits a consumer's unpaid medical debt amount, account number, and home address electronically to a third-party commercial mailing vendor so the vendor can automatically print, stuff, and mail a collection letter, allegedly violating a federal ban on unauthorized third-party debt disclosures. The consumer sues for statutory damages, alleging emotional distress and invasion of privacy, though no human at the vendor is proven to have read the file. This is a classic post-\textit{TransUnion} fence-sitter: plaintiffs argue that disclosing private debt data to an outside corporation closely parallels the privacy tort of public disclosure of private facts (and Thomas's dissent invoked common-law publication to stenographers and telegraph operators), whereas defendants invoke \textit{TransUnion}'s footnote 6 to argue that automated processing by a ministerial print vendor is neither ``public'' disclosure nor ``read'' by a human, and that subjective distress alone is insufficient (see \href{https://www.supremecourt.gov/DocketPDF/23/23A952/308036/20240419131911550_Lower\%20Court\%20Opinions.pdf}{\textit{Barclift v. Keystone Credit Services, LLC}}, 93 F.4th 136 (3d Cir. 2024)).
\end{enumerate}

\section{Critique}

\textbf{Originalist and Textualist Critique:} A striking feature of \textit{TransUnion} is that its sharpest originalist attack came from Justice Thomas in dissent, echoed by prominent originalist and historical scholarship. Critics point out that the phrase ``injury in fact'' appears nowhere in Article III and was invented in 1970 (\textit{Data Processing}), whereas Founding-era courts routinely adjudicated suits for violations of private statutory rights---such as the First Congress's copyright, patent, and penal statutes---without requiring proof of actual or common-law damage (\href{https://journals.librarypublishing.arizona.edu/arizlrev/article/10580/galley/9221/download/}{Arizona Law Review}; \href{https://scholarship.law.ufl.edu/cgi/viewcontent.cgi?article=4189&context=flr}{Florida Law Review}). Scholars further note a tension between \textit{TransUnion} and the Court's own historical methodology in \textit{Vermont Agency of Natural Resources v. United States ex rel. Stevens}, 529 U.S. 765 (2000), and \textit{SEC v. Jarkesy}, 603 U.S. 109 (2024): if Founding-era penal statutes authorizing statutory penalties without proof of actual harm were traditional legal actions belonging exclusively in Article III courts (including under \S~9 of the Judiciary Act of 1789), then \textit{TransUnion}'s historical rationale rests on an incomplete record (\href{https://journals.librarypublishing.arizona.edu/arizlrev/article/10580/galley/9221/download/}{Arizona Law Review}). Commentators have also observed that the majority's ``close historical analogue'' test operates more as a Burkean presumption against legislative novelty than genuine originalism, inviting judges to freeze Congress's Article I remedial authority at 18th-century common-law baselines (\href{https://harvardlawreview.org/print/vol-137/the-presumption-against-novelty-in-the-roberts-courts-separation-of-powers-case-law/}{Harvard Law Review (2024)}; \href{https://lawreview.gmu.edu/print__issues/the-court-continues-to-confuse-standing-the-pitfalls-of-faux-article-iii-originalism/}{George Mason Law Review}).

\textbf{Progressive and Doctrinal Critique:} Progressive and doctrinal scholars criticize \textit{TransUnion} on three main grounds. First, they argue that the ``close relationship'' test is doctrinally indeterminate: the majority offered no principled metric explaining why common-law defamation's \textit{falsity} element was dispensable (allowing ``potential match'' alerts to qualify) while its \textit{publication} element was indispensable, leaving lower courts ``at sea'' and producing circuit splits in data-privacy, wiretap, and consumer-protection litigation (\href{https://harvardlawreview.org/print/vol-136/standing-in-the-way-the-courts-escalating-interference-in-federal-policymaking/}{Harvard Law Review (2023)}; \href{https://btlj.org/wp-content/uploads/2026/08/Hua_Ready-ADA-Compliant_08-11-26.pdf}{Berkeley Technology Law Journal}). Second, critics contend that the majority misapplied separation-of-powers principles: unlike \textit{Lujan}, which involved a citizen suit challenging Executive Branch inaction, the FCRA creates individualized private rights against private corporations, posing no threat to the President's Article II Take Care duties (\href{https://harvardlawreview.org/print/vol-135/transunion-v-ramirez/}{Harvard Law Review (2021)}). Third, as Justice Thomas warned and the \textit{Harvard Law Review} emphasized, barring federal standing for absent class members undermines the Class Action Fairness Act of 2005 (CAFA) by pushing nationwide federal statutory class actions out of federal court and into state courts that are not bound by Article III and from which defendants cannot remove (\href{https://harvardlawreview.org/print/vol-135/transunion-v-ramirez/}{Harvard Law Review (2021)}).

\textbf{Underlying Competing Values:} Beneath the doctrinal debate lies a clash between two foundational constitutional values. On one side is \textit{limited judicial power and anti-abusive litigation}: the majority seeks to prevent Congress from deputizing unharmed private ``bounty hunters'' and class-action counsel to extract massive statutory windfalls from private businesses, preserving prosecutorial discretion for politically accountable Executive officials and reserving federal dockets for tangible, real-world disputes. On the other side is \textit{democratic legislative supremacy and modern prophylactic deterrence}: in an era of automated data brokers, algorithmic scoring, and mass digital surveillance, Congress often chooses statutory damages and private enforcement precisely because modern informational harms are diffuse, hard to quantify, and too vast for under-resourced federal agencies to police alone. \textit{TransUnion} places the federal judiciary---rather than Congress---as the ultimate gatekeeper of which 21st-century harms are ``real'' enough to enter federal court (\href{https://harvardlawreview.org/print/vol-136/standing-in-the-way-the-courts-escalating-interference-in-federal-policymaking/}{Harvard Law Review (2023)}; \href{https://washulawreview.org/the-evolution-of-standing-and-the-need-for-foundational-realism-a-historical-inquiry-from-hayburn-to-harvard/}{Washington University Law Review}).

\section{Key Quotations}

\begin{quote}
``To have Article III standing to sue in federal court, plaintiffs must demonstrate, among other things, that they suffered a concrete harm. No concrete harm, no standing.''
\end{quote}
\noindent --- Justice Kavanaugh, Opinion of the Court, 594 U.S. 413, 417 (2021) (slip op. at 1) (\href{https://www.supremecourt.gov/opinions/20pdf/20-297_4g25.pdf}{Slip Opinion}; \href{https://supreme.justia.com/cases/federal/us/594/20-297/}{Justia}).

\begin{quote}
``Congress may enact legal prohibitions and obligations. And Congress may create causes of action for plaintiffs to sue defendants who violate those legal prohibitions or obligations. But under Article III, an injury in law is not an injury in fact. Only those plaintiffs who have been concretely harmed by a defendant's statutory violation may sue that private defendant over that violation in federal court.''
\end{quote}
\noindent --- Justice Kavanaugh, Opinion of the Court, 594 U.S. 413, 427 (2021) (slip op. at 11) (\href{https://www.supremecourt.gov/opinions/20pdf/20-297_4g25.pdf}{Slip Opinion}; \href{https://www.law.cornell.edu/supremecourt/text/20-297}{Cornell LII}).

\begin{quote}
``In cases such as these where allegedly inaccurate or misleading information sits in a company database, the plaintiffs' harm is roughly the same, legally speaking, as if someone wrote a defamatory letter and then stored it in her desk drawer. A letter that is not sent does not harm anyone, no matter how insulting the letter is.''
\end{quote}
\noindent --- Justice Kavanaugh, Opinion of the Court, 594 U.S. 413, 434 (2021) (slip op. at 18--19) (\href{https://www.supremecourt.gov/opinions/20pdf/20-297_4g25.pdf}{Slip Opinion}; \href{https://supreme.justia.com/cases/federal/us/594/20-297/}{Justia}).

\begin{quote}
``Never before has this Court declared that legal injury is inherently insufficient to support standing. And never before has this Court declared that legislatures are constitutionally precluded from creating legal rights enforceable in federal court if those rights deviate too far from their common-law roots. According to the majority, courts alone have the power to sift and weigh harms to decide whether they merit the Federal Judiciary's attention. In the name of protecting the separation of powers, ante, at 422--423, 429, this Court has relieved the legislature of its power to create and define rights.''
\end{quote}
\noindent --- Justice Thomas, dissenting, 594 U.S. 413, 453--54 (2021) (footnote omitted) (\href{https://www.supremecourt.gov/opinions/20pdf/20-297_4g25.pdf}{Slip Opinion}; \href{https://supreme.justia.com/cases/federal/us/594/20-297/}{Justia}).

\begin{quote}
``The Court here transforms standing law from a doctrine of judicial modesty into a tool of judicial aggrandizement. It holds, for the first time, that a specific class of plaintiffs whom Congress allowed to bring a lawsuit cannot do so under Article III.''
\end{quote}
\noindent --- Justice Kagan, dissenting, 594 U.S. 413 (2021) (141 S. Ct. 2190, 2225; slip op. at 1) (\href{https://www.supremecourt.gov/opinions/20pdf/20-297_4g25.pdf}{Slip Opinion}; \href{https://www.law.cornell.edu/supremecourt/text/20-297}{Cornell LII}).

\end{document}
